\documentclass[tikz,border=0pt]{standalone}
\def\FigureLanguage{bi}
\def\FigurePrimary{zh}
% Put before academic-preamble.tex; override these commands before input if needed.
\providecommand{\FigureLanguage}{zh}
\providecommand{\FigurePrimary}{zh}
\newif\ifFigureEnglishOnly
\newif\ifFigureBilingual
\newif\ifFigureEnglishFirst
\def\FigureLangEnglish{en}
\def\FigureLangChinese{zh}
\def\FigureLangBilingual{bi}
\edef\FigureSelectedLanguage{\FigureLanguage}
\edef\FigureSelectedPrimary{\FigurePrimary}
\ifx\FigureSelectedLanguage\FigureLangEnglish
  \FigureEnglishOnlytrue
\else\ifx\FigureSelectedLanguage\FigureLangBilingual
  \FigureBilingualtrue
\else\ifx\FigureSelectedLanguage\FigureLangChinese
\else
  \PackageError{paperfigure}{Unknown FigureLanguage: \FigureLanguage}{Use zh, en, or bi.}
\fi\fi\fi
\ifx\FigureSelectedPrimary\FigureLangEnglish
  \FigureEnglishFirsttrue
\else\ifx\FigureSelectedPrimary\FigureLangChinese
\else
  \PackageError{paperfigure}{Unknown FigurePrimary: \FigurePrimary}{Use zh or en.}
\fi\fi
\typeout{FIGURE-LANGUAGE: \FigureLanguage}
\typeout{FIGURE-PRIMARY: \FigurePrimary}
% Ordinary text must be TeX-escaped; shared math is trusted TeX.
\newcommand{\FigureLabelContent}[3]{%
  \ifFigureBilingual
    \ifFigureEnglishFirst #2\\[.25ex]#1\else #1\\[.25ex]#2\fi
  \else
    \ifFigureEnglishOnly #2\else #1\fi
  \fi
  \if\relax\detokenize{#3}\relax\else\\[.25ex]$#3$\fi}
% Optional width enables genuine paragraph wrapping, without reducing type size.
\newcommand{\FigureLabel}[4][]{%
  \if\relax\detokenize{#1}\relax
    \shortstack[c]{\FigureLabelContent{#2}{#3}{#4}}%
  \else
    \parbox[c]{#1}{\centering\FigureLabelContent{#2}{#3}{#4}\par}%
  \fi}

\def\FigureCJKFont{SimSun}
\def\FigureLatinFont{Times New Roman}
\def\FigureMathFont{TeX Gyre Termes Math}
% Input after documentclass, before begin{document}; XeLaTeX/LuaLaTeX.
% Authors may define the Figure*Font commands before input to override defaults.
% English-only documents can set FigureLanguage=en without loading CJK support.
\providecommand{\FigureLanguage}{zh}
\providecommand{\FigurePrimary}{zh}
\makeatletter
\@ifundefined{ifFigureEnglishOnly}{%
  \newif\ifFigureEnglishOnly
  \def\FigureFontEnglishMode{en}%
  \edef\FigureFontSelectedMode{\FigureLanguage}%
  \ifx\FigureFontSelectedMode\FigureFontEnglishMode\FigureEnglishOnlytrue\fi
}{}
\makeatother
\providecommand{\FigureCJKFont}{SimSun}
\providecommand{\FigureLatinFont}{Times New Roman}
\providecommand{\FigureMathFont}{TeX Gyre Termes Math}
\ifFigureEnglishOnly
  \usepackage{fontspec}
  \typeout{FIGURE-CJK-FONT: not requested (English-only)}
\else
  \usepackage[UTF8,fontset=none]{ctex}
\fi
\usepackage{amsmath}
\usepackage{unicode-math}
\ifFigureEnglishOnly\else
\IfFontExistsTF{\FigureCJKFont}
  {\setCJKmainfont{\FigureCJKFont}\typeout{FIGURE-CJK-FONT: \FigureCJKFont}}
  {\setCJKmainfont{FandolSong-Regular}
   \typeout{FIGURE-CJK-FONT: FandolSong-Regular (fallback)}
   \PackageWarningNoLine{paperfigure}{Chinese font fallback: FandolSong-Regular}}
\fi
\IfFontExistsTF{\FigureLatinFont}
  {\setmainfont{\FigureLatinFont}\typeout{FIGURE-LATIN-FONT: \FigureLatinFont}}
  {\setmainfont{Latin Modern Roman}
   \typeout{FIGURE-LATIN-FONT: Latin Modern Roman (fallback)}
   \PackageWarningNoLine{paperfigure}{Latin font fallback: Latin Modern Roman}}
\IfFontExistsTF{\FigureMathFont}
  {\setmathfont{\FigureMathFont}\typeout{FIGURE-MATH-FONT: \FigureMathFont}}
  {\setmathfont{latinmodern-math.otf}
   \typeout{FIGURE-MATH-FONT: Latin Modern Math (fallback)}
   \PackageWarningNoLine{paperfigure}{Requested math font unavailable; using Latin Modern Math}}
% For individual bold letters/Greek symbols, not a replacement for bold equations.
% symbf follows bold-style; use symbfup/symbfit for explicit upright/italic.
\providecommand{\FigureBoldSymbol}[1]{\symbf{#1}}

\usepackage{tikz}
\usetikzlibrary{arrows.meta}
\begin{document}
\fontsize{9.5}{11.875}\selectfont
\ifFigureBilingual
  \def\H{102}\def\Top{13}\def\TOne{24}\def\TTwo{45}\def\Bottom{58}\def\Note{68}\def\Rule{78}\def\Eq{89}
\else
  \def\H{92}\def\Top{11}\def\TOne{21}\def\TTwo{40}\def\Bottom{53}\def\Note{61}\def\Rule{69}\def\Eq{81}
\fi
\begin{tikzpicture}[x=1mm,y=-1mm,
 every node/.style={font=\fontsize{9.5}{11.875}\selectfont,inner sep=0pt},
 link/.style={line width=.65pt},
 flow/.style={-{Stealth[length=2mm,width=1.4mm]},line width=.65pt}]
\path[use as bounding box] (0,0) rectangle (148,\H);
\node at (74,5) {\FigureLabel{并联导热的等效热阻网络}{Equivalent network for parallel heat conduction}{}};
% Two reservoirs are physical terminals, not start/end workflow nodes.
\draw[link] (24,\Top+6) -- (24,\Bottom);
\draw[link] (124,\Top+6) -- (124,\Bottom);
\node at (13,\TOne+10) {$T_{\mathrm h}$};
\node at (136,\TOne+10) {$T_{\mathrm c}$};
\foreach \y/\r/\q in {\TOne/1/1,\TTwo/2/2} {
  \draw[link] (24,\y) -- (55,\y);
  \draw[link,fill=black!7] (55,\y-3) rectangle (93,\y+3);
  \draw[link] (93,\y) -- (124,\y);
  \node at (74,\y) {$R_{\r}$};
  \fill (24,\y) circle[radius=.7mm];
  \fill (124,\y) circle[radius=.7mm];
  \draw[flow] (31,\y-4) -- (46,\y-4);
  \node at (39,\y-8) {$\dot Q_{\q}$};
}
\node at (74,\Bottom) {\FigureLabel{两条支路连接相同温度端点}{Both branches connect the same temperature terminals}{}};
\node at (74,\Note) {\FigureLabel{稳态；理想端点；示意，不按比例}{Steady state; ideal terminals; schematic, not to scale}{}};
\draw[black!50,line width=.45pt] (6,\Rule) -- (142,\Rule);
\node at (40,\Eq) {$\displaystyle \dot Q_i=\frac{T_{\mathrm h}-T_{\mathrm c}}{R_i},\quad R_i>0$};
\node at (109,\Eq) {$\displaystyle \dot Q=\dot Q_1+\dot Q_2$};
\end{tikzpicture}
\end{document}
